Israel, which allocates 8.8\% of its GDP to defense~\citep{sipri2025milex}, invests heavily in the development of autonomous weapon systems (AWS), with their Prime Minister announcing a \$100 billion investment in homegrown arms production last December~\citep{netanyahu2025arms}. Israel's most well-known homegrown AWS is the Iron Dome air defense system, which intercepts incoming projectiles in real time.  Building on this foundation in artificial intelligence (AI)-assisted defense technology, Israel has also expanded into offensive technology, such as the `Habsora' targeting system.

Habsora (`The Gospel') is an AI-targeting system used by the Israeli Defense Force (IDF), developed primarily by Unit 8200 within the Military Intelligence Directorate. First deployed during the May 2021 Operation `Guardian of the Walls', it generates target recommendations using a multimodal intelligence pipeline that encompasses signal intelligence, telephone metadata, satellite and drone imagery, social graph analysis, and life-pattern modelling \cite{abraham_2024, dwoskin_2024}. The system is an ensemble of multiple specialised algorithms fused together \cite{dwoskin_2024}. During the 2023-2024 Gaza conflict, Habsora was reported to generate over one hundred target recommendations per day, which in previously would require approximately twelve months of manual analysis \cite{abraham_2024, ynet_2024}.

Habsora's AI capability is the probabilistic classification and prioritisation of locations and individuals as military objectives. It assigns confidence scores derived from large-scale pattern matching across a shared intelligence pool, rather than positive identification in the traditional intelligence analysis \cite{ootlaws}. Military commanders authorise each recommendation, yet the review window documented for each target was approximately only twenty seconds \cite{abraham_2024}.

This creates a stakeholder dynamic with both ethical and legal significance. Its output influencing both IDF commands, who operate under high cognitive load and time pressure; Palestinian civilians in the target environment bearing the lethal consequences and the broader international community, which depends on International Humanitarian Law (IHL) compliance for the legitimacy of armed conflict. \citet{mhcof} offer a framework for this tension. They propose a two-axis model of Meaningful Human Control (MHC) which distinguishes between high automation and genuine human control. Both require adequate information and real freedom of action. The central argument of this report is that Habsora structurally undermines the second condition, and that this gap carries cascading implications across the technical, sociotechnical, and governance layers identified by \citet{info12090385} in their Comprehensive Human Oversight Framework (CHOF).